%% VII. Conclusion and future work. Input by access.tex.
\section{Conclusion}
\label{sec:conclusion}
Under two locked preregistrations and eight verification gates, a linear-Gaussian state-space model of meditation dynamics, derived from the BQI framework, was fitted to experience-sampling EEG from two laboratories. On the four research questions: (RQ1) the precision predicted by an exact Fisher-information design analysis before the EEG was seen was achieved on ds001787 and missed by 28\,\% on OSF 3uszv (8.5 against 6.7 percent for the DMN rate), well inside the factor of two with which the bound was validated in simulation, although the observed-information and profile errors expose a misspecified alpha-power channel; (RQ2) neither EEG proxy carries information about the reports beyond their own serial dependence; (RQ3) of the five predictions of BQI for the estimated quantities (Section~\ref{sec:theory-pred}), only the positive sign of $\alpha$ is reproduced, with a determined sign on ds001787 alone, while the time constants, the coupling, the sign of $\beta$ and the group difference are not; and (RQ4) the rate constants, estimated from uncalibrated proxies, are properties of the feature in the recording rather than of a brain process, so they do not transfer between laboratories. This is the first empirical support for using such a design analysis to plan experience-sampling studies, and a demonstration of what it cannot promise. Design analyses of dynamical models of mental states should therefore be read as statements about precision, be accompanied by observed-information and profile errors, and be preceded by a calibration of the features that is shown to hold across recordings.

Three lines of future work follow. First, a calibration study that records EEG together with fMRI and a complexity index, as BQI proposes \cite[Sec.~14]{BQI}, would establish which EEG feature, if any, tracks DMN activity and integration well enough for the design analysis to apply. Second, a hierarchical state-space model with subject-level rate constants would address the large between-subject heterogeneity and change the design map quantitatively. Third, a non-linear or ordinal read-out would match the bounded and, on one dataset, bimodal ratings, at the price of approximate filters and simulation-based Fisher information.
